% -------------
\section{Conclusion and Future Work}
\label{sec:conclusion}

We introduced \ourdataset, a symbolic benchmark that pairs each engineering problem with an exact
gold reasoning trace, and an evaluator that checks responses against this trace, calling an LLM
judge only where matching cannot decide.
The strongest LLMs were hard to tell apart by their final answers, yet their derivations differed in
how closely they followed the gold route and how often their displayed arithmetic was flagged.
Most failures were calculation slips rather than conceptual errors, which were rarer on easy
problems.
Reasoning raised the final-answer accuracy of every model that used it, and the order of the models
held under rewording, a changed answer tolerance, and repeated decoding.
These results suggest that, near the ceiling, progress in engineering reasoning is better measured
in the derivation than in the final answer.
Future work will check the principles and assumptions a solution applies, not only its numbers,
since a misstated rule behind a correct result escapes every numerical test, and will turn the gold
reasoning traces into feedback for training and selecting LLMs.
Moreover, extending the approach to open-ended design, real engineering artifacts, and tool-based
workflows would bring evaluation closer to engineering practice.
